\documentclass[11pt]{article}
\usepackage[T1]{fontenc}
\usepackage{lmodern,amsmath,amssymb,amsthm,booktabs,geometry,xurl}
\usepackage[hidelinks]{hyperref}
\geometry{margin=1in}
\newtheorem{theorem}{Theorem}
\newtheorem{lemma}[theorem]{Lemma}
\newcommand{\file}[1]{\texttt{\nolinkurl{#1}}}
\newcommand{\PA}{\mathrm{PA}}
\setlength{\emergencystretch}{2em}
\title{A validated regrouping of PA196\\and its arithmetic transfer}
\author{Technical note}
\date{13 November 2026}
\begin{document}
\maketitle

\begin{abstract}
This note records a new 196-working-state layout of the unchanged
PA196 register program.  Regrouping the binary program tree saves
one reader state, but its longest zero run becomes seven bits, which
requires an eighth-zero return test and consumes that saving.  The
result passes all $4{,}959$ uniform symbolic instruction contracts
at minimum gap $48$, and reaches the exact baseline clean frame
from blank tape after $59{,}164$ steps.  It is a verified intermediate
construction, not a machine below 196 states.  We explain what the
symbolic certificates establish, how their arithmetic proofs transfer
the halting statement, and which steps have not been kernel-checked.
\end{abstract}

\section{Pinned artifact and state convention}

The artifact is \file{data/variant196.native.tm}.  It has 196
working controls and 392 rows, starts at native control $74$ on an
all-zero two-way infinite binary tape, and uses the separate halt
label $254$.  The native labels are sparse; they do not determine
the state count by their maximum.  Its SHA-256 digest is
\begin{quote}\small\nolinkurl{f6594edc95efed02f26b88f71c9b1f7699965a258739836f24e6b641887b0cd9}.
\end{quote}
The matching source is \file{data/variant196_source.json}.
Its SHA-256 digest is
\begin{quote}\small\nolinkurl{f2fd582615c47a680c6475235139d8026bd4126b03e788e7d18c76c1ff45c49e}.
\end{quote}
A finite bijection can rename the 196 working controls to
$0,\ldots,195$, send native start $74$ to dense start $0$, and
send the separate halt to $196$.  Checking the bijection against
all rows proves literal trace equivalence; no state is removed by
that operation.

The baseline is James Chang's PA196 at commit
\href{https://github.com/jamesyc/BusyBeaver/tree/407cdc356f6b586635e71ac92be394255f4bfea4}{\nolinkurl{407cdc356f6b586635e71ac92be394255f4bfea4}}.
The local baseline native table has SHA-256
\begin{quote}\small\nolinkurl{454254a42a9f9096b72f73ef491107713002ca32f2a51f165eee7ca299283215}.
\end{quote}
This native-table digest differs from that of the upstream dense
\file{data/pa196.tm} because the labels and serialization differ.

\section{What changed}

The $3{,}812$ primitive source instructions retain their ordinal
order, opcodes, register arguments and resolved successors.  The
axiom command emits the same formulas in the same order.  Only
binary parenthesization and physical sharing of the reader states
change.  Two independent successor resolvers, one using binary
address arithmetic and one using subtree intervals, give the same
control-flow graph as the baseline, including both outcomes of
conditional decrement.

The regrouped source has 158 distinct internal binary nodes, one
fewer than the baseline's 159.  It has 23 distinct instruction leaves;
thus the runtime checker's reported total of 181 distinct DAG nodes
includes leaves and is consistent with the internal-node count.
Six selected context folds in the regrouped AX region share
complementary reader occurrences; inherited sharing outside this
region is retained.  These are contextual equivalences, not ordinary
merges of identical transition rows.  Their correctness is included
in the complete instruction contracts below.

The maximum source-path depth is 36 and the maximum consecutive
zero run is 7.  A seven-zero return scanner can mistake such an
internal path run for the blank exterior and resume at the wrong
path position.  Hence the nominal 195-state layout is not accepted.
The repaired table inserts the following eighth test, in native
labels:
\[
 \begin{array}{c|cc}
       &\text{read }0&\text{read }1\\ \hline
 34&(0,L,2000)&(1,L,2)\\
 2000&(0,L,75)&(1,L,2).
 \end{array}
\]
The preceding zero-count states are
$2,3,5,9,15,23,34$.  On a one they reset leftward to state $2$.
The extra state $2000$ distinguishes a path run of seven zeros from
the longer blank run used for returning to the root.  This costs
one working state, bringing the repaired table back to 196.
A static maximum-zero count is useful for the design, but the full
physical contracts, rather than that count alone, validate the repair.

\section{The checked results}

The receipts \file{certificates/variant196/runtime.json} and
\file{certificates/variant196/startup.json} record:
\begin{center}
\begin{tabular}{@{}lr@{}}
\toprule
Quantity & Result\\
\midrule
Working states / table rows & 196 / 392\\
Source sites / decrement sites & 3,812 / 1,147\\
Uniform instruction cases & 4,959\\
Rows covered by the symbolic contracts & 392\\
Minimum permitted header gap & 48\\
Maximum source-path depth / zero run & 36 / 7\\
Symbolic inference steps & 1,177,411\\
Smallest constant term in a contract duration & 3\\
Blank-start steps to the baseline clean frame & 59,164\\
\bottomrule
\end{tabular}
\end{center}
The $4{,}959$ cases consist of one case per source site plus the
second outcome at each decrement site.  Every case closes for nine
independent arbitrary natural register values, an arbitrary first
spare value, an arbitrary finite remaining spare list, and every
gap $48+g$ with $g\in\mathbb N$.  This is not a finite sampling of
small register contents.

Startup is a separate exact finite simulation from the blank tape.
After $59{,}164$ steps the new table has control $84$, head position
$-2$, and exactly the same 66 one-cells as the baseline after its
$58{,}154$-step startup.  The first divergence in their literal runs
is at step 29, so the startup claim is equality of the final clean
frame, not equality of the two complete traces.  The initial source
site $25$ retains path \texttt{100000}, so the common physical frame
represents the same initial source configuration.  Its gap is 48,
within the uniformly checked range.

\section{Soundness of the symbolic certificate rules}

A symbolic tape contains finite known bits, runs whose lengths are
affine expressions in independent natural parameters, a finite list
of spare register blocks, and the infinite blank exterior.  The
checker accepts a literal read from a variable run only when its
length is known to be positive, or handles the whole run by an
explicit self-loop induction.  A proof that needs to inspect an
unrepresented or possibly empty region is rejected.

The elementary proof rules are as follows.
\begin{enumerate}
\item A literal transition is the corresponding row of the supplied
machine table and changes precisely one scanned cell and the head.
\item A unary self-loop traverses a whole run.  If its row is
$(q,b,w,D,q)$, induction on $n$ proves its exact action on $b^n$,
including $n=0$, its duration $n$, and its displacement $nD$.
\item The four generic one-register traversals are rechecked from
the actual table.  They act on
$B(v)=0\,1^{v+1}$ or $C(v)=1\,0\,1^v$, have duration $v+2$,
and retain opaque contexts on both sides.  Empty/cons induction
extends each one-block theorem to every finite spare list.
\item Adjacent equal-bit runs combine by addition of their lengths.
Reversal, the identities $C(A)1=1B(A)$ for lists of blocks and their
reversed counterparts, and absorption of finite zero runs into the
blank exterior preserve the represented tape.  These are direct
word identities, proved by finite-list induction where necessary.
\end{enumerate}
Composition of these steps proves every complete contract.  The
reported durations have nonnegative affine coefficients and positive
constant terms, ensuring genuine positive physical progress.
The target configurations show exact register effects and successor
paths, with nondecreasing gap and the required spare-file invariant.
The four generic block lemmas and all site contracts are rederived
from the new table; no unchanged name of an engine state is used
as an unverified theorem.

All the above proofs are ordinary inductions on finite encodings.
They are formalizable in PA, including for nonstandard parameter
values in a PA model.  This mathematical soundness argument, together
with successful certificate replay, justifies the uniform contracts.
The Python checker is part of the trusted proof boundary.  Its
successful execution is not being described as a Lean or Rocq
kernel check.

% Report-ready fragment, requiring amsmath and the existing \file macro.
\paragraph{An explicit bound for the arithmetic trace construction.}
The affine duration vectors give a concrete bound on the PA witness
construction.  At every live boundary after $n$ primitive source steps
from the checked clean seed, the header-gap excess over $48$ is at most
$n$, the sum of the nine
register values is at most $n+2$, the first spare remains $1$, and the
width of the remaining spare list is at most $75+3n$.  The initial spare
file has 26 entries of value 1, and each block $B(1)=011$ has width 3.
Every increment appends one such block and raises one register by one;
every positive decrement increases the gap by one and lowers one
register by one; the other branches preserve these quantities.
All source-path increments and cuts satisfy the checked no-root-overflow
condition.  Without that condition, an additional header-displacement
term would have to be included; the present estimate does not silently
drop such a term.

Write a checked duration vector as
$(a_0,a_g,a_1,\ldots,a_9,a_s,a_w)$ and put
$m=\max(a_1,\ldots,a_9)$.  Nonnegative coefficients and the preceding
invariants bound this primitive's duration by
\[
  (a_0+2m+a_s+75a_w)+(a_g+m+3a_w)n.
\]
In a positive-decrement case, the tested register parameter represents
its value minus one, so this estimate remains valid.  In a zero-decrement
case its unused parameter has coefficient zero, which is checked
explicitly.  The standalone \file{code/check_trace_bound.py} checks all
$4{,}959$ vectors, their complete case census, the full canonical seed,
and every relevant no-overflow source path.  The maximum constant is
$700$ and the maximum slope is $18$.  Hence a trace implementing a live
prefix of $T$ primitive source steps, or one whose final instruction
halts, has length, including blank startup, at most
\[
 \boxed{59{,}164+700T+9T(T-1)}.
\]
This is the exact sum of $700+18n$ for $0\leq n<T$, plus startup.
It gives an elementary PA-provable bound for producing the literal
trace witness.  The conclusion relies on the established uniform runtime
contracts; checking these numerical inequalities is not a replacement
for their symbolic induction proofs and does not imply nonhalting by
bounded simulation.

\section{Arithmetic transfer}

Let $R$ denote this repaired table, $M_{196}$ the baseline, and $S$
their identical ordinal source program.  Let $H_R,H_{196},H_S$ be
the respective exact coded halting sentences.

\begin{theorem}
The source identity, startup trace and uniform positive contracts,
with the preceding PA-formalizable soundness proof, give
\[
 \PA\vdash H_R\leftrightarrow H_S\leftrightarrow H_{196}.
\]
Consequently the published internal consistency implication transfers:
\[
 \PA\vdash\neg H_R\to\operatorname{Con}(\PA).
\]
With the direct uniform syntactic PA proof translation for the native
calculus, described in the accompanying soundness note, the full
paper-level conclusion is
\[
 \PA\vdash\neg H_R\leftrightarrow\operatorname{Con}(\PA).
\]
\end{theorem}
\begin{proof}
A finite source halt gives a literal halt by prepending startup and
concatenating its instruction contracts.  Conversely, if a literal
halt occurs at time $T$ but the source does not halt, PA constructs
$T+1$ source steps and concatenates their positive literal segments.
The result is a genuine literal trace longer than $T$ with a working
endpoint.  Determinism identifies its configuration at time $T$
with the alleged halt, which has no outgoing transition.  This
contradiction proves the converse.  Both arguments use PA induction
on coded finite sequences, so they handle nonstandard lengths.
Apply them to both tables to obtain the first equivalences.

The upstream development proves the internal direction
$\neg H_{196}\to\operatorname{Con}(\PA)$, which transfers immediately.
For the reverse direction, the separate syntactic translation maps
membership to binary digits and compiles every guarded separation
instance into a PA proof using finite bit comprehension.  It retains
proof witnesses beside all live source formulas.  Under the internal
hypothesis $\operatorname{Con}(\PA)$, Shift cannot reach the false
code $1$, so all command boundaries remain live.  The positive
command/source and source/tape contracts then give arbitrarily long
live literal traces and exclude halting by the same determinism
argument.
\end{proof}

\end{document}
